% =====================================================================
\chapter{Soluciones}
\label{ap:soluciones}
% =====================================================================

Las remisiones del tipo «Módulo 0.5, §5.2» indican el módulo de la fase~0 y la sección de sus apuntes (archivos P01 a P09) tal como aparecen en el PDF. Cuando se pide un valor decimal, los resultados finales se dan con cuatro cifras significativas. Al final de cada solución figuran, entre corchetes, los criterios de corrección.

% ---------------------------------------------------------------------
\section{Parte A: clave de respuestas}

\begin{table}[h]
\centering\footnotesize
\caption{Clave de la Parte~A.}\label{tab:clave}
\begin{tabularx}{\textwidth}{@{}c c >{\raggedright\arraybackslash}X l@{}}
\toprule
N.º & Resp. & Justificación & Repasar \\
\midrule
1 & (a) & $\neg\forall\equiv\exists\neg$ y $\neg\exists\equiv\forall\neg$; la negación de $<$ es $\ge$. & Módulo 0.1, §2.4 \\
2 & (c) & El contrarrecíproco de $P\to Q$ es $\neg Q\to\neg P$; (a) es el recíproco y (b), el inverso. & Módulo 0.1, §1.2 \\
3 & (d) & $2a+1=2b+1\Rightarrow a=b$; la imagen solo contiene impares, así que $2$ no tiene preimagen. & Módulo 0.2, §3.1 \\
4 & (a) & $f(a)=f(a')\Rightarrow g(f(a))=g(f(a'))\Rightarrow a=a'$; $g$ puede no ser inyectiva. & Módulo 0.2, §4.2 \\
5 & (b) & La suma de los $n$ primeros impares es $n^{2}=20^{2}=400$. & Módulo 0.3, §4.3 \\
6 & (c) & $\log_2(x^{2}/y)=2\log_2x-\log_2y=10-3=7$. & Módulo 0.3, §5.2 \\
7 & (d) & $5=\texttt{00000101}$; se invierten los bits (\texttt{11111010}) y se suma~1. & Módulo 0.3, §6.3 \\
8 & (b) & $\cos(45^{\circ}+30^{\circ})=\frac{\sqrt2}{2}\frac{\sqrt3}{2}-\frac{\sqrt2}{2}\frac12=\frac{\sqrt6-\sqrt2}{4}=0{,}2588$. & Módulo 0.4, §2.2 \\
9 & (c) & Suma geométrica de razón $\omega=e^{2\pi i/n}\neq1$: $(1-\omega^{n})/(1-\omega)=0$. & Módulo 0.4, §6.5 \\
10 & (a) & $1-\cos x=2\sin^{2}(x/2)$ y $\sin(x/2)/(x/2)\to1$: el límite es $2\cdot\frac14=\frac12$. & Módulo 0.5, §2.4 \\
11 & (d) & $\sigma'=e^{-x}/(1+e^{-x})^{2}=\sigma\cdot e^{-x}/(1+e^{-x})=\sigma(1-\sigma)$. & Módulo 0.5, §4.3 \\
12 & (b) & $\sin0=0$, $\sin'0=1$, $\sin''0=0$, $\sin'''0=-1$: $x-x^{3}/3!$. & Módulo 0.5, §6.1 \\
13 & (c) & Con $u=x^{2}$: $\frac12\int_0^1e^{u}\,\dd u=\frac{e-1}{2}=0{,}8591$. & Módulo 0.6, §3.1 \\
14 & (b) & Serie $p$ con $p=1{,}01>1$; (a) y (c) son series $p$ con $p\le1$ y en (d) el término no tiende a~0. & Módulo 0.6, §5.3 \\
15 & (a) & $\nabla f=(2xy,\ x^{2}+e^{y})$, que en $(1,0)$ vale $(0,\,1+1)=(0,2)$. & Módulo 0.7, §4.4 \\
16 & (c) & $\det H=4-9=-5<0$: los autovalores tienen signos opuestos. & Módulo 0.7, §5.2 \\
17 & (d) & Permutaciones con repetición: $6!/(3!\,2!\,1!)=60$. & Módulo 0.8, §2.2 \\
18 & (c) & $\frac{0{,}01\cdot0{,}99}{0{,}01\cdot0{,}99+0{,}99\cdot0{,}05}=\frac{0{,}0099}{0{,}0594}=0{,}1667$. & Módulo 0.8, §4.3 \\
19 & (d) & $\mathbb E[\widehat\sigma^{2}]=\frac{n-1}{n}\sigma^{2}=\frac45\cdot10=8$ (estimador sesgado). & Módulo 0.9, §2.5 \\
20 & (b) & $50\pm1{,}96\cdot10/\sqrt{100}=50\pm1{,}96$; (a) usa $1{,}645$ (nivel 90\,\%). & Módulo 0.9, §4.3 \\
\bottomrule
\end{tabularx}
\end{table}

% ---------------------------------------------------------------------
\section{Parte B: soluciones de los problemas}

% ------------------------------------------------------------- B.1
\begin{solucion}{eje:logica}
\datos proposiciones $p,q,r$; la definición de continuidad en~$a$; $n\in\Z$; la sucesión de Fibonacci con $F_0=0$, $F_1=1$, $F_{n+2}=F_{n+1}+F_n$.
\planteamiento (a) se comparan columnas de una tabla de verdad (Módulo 0.1, §1.3 y §1.4); (b) se aplica la negación de cuantificadores (Módulo 0.1, §2.4) y $\neg(P\to Q)\equiv P\wedge\neg Q$; (c) se demuestra el contrarrecíproco (Módulo 0.1, §3.3); (d) se usa inducción simple sobre la conjunción de las dos identidades (Módulo 0.1, §5.1).
\pasos

\noindent\textbf{(a)} Basta comparar las columnas de las dos fórmulas en las ocho filas:
\begin{center}\small
\begin{tabular}{@{}ccc|cc|c|c|c@{}}
\toprule
$p$ & $q$ & $r$ & $p\to q$ & $p\to r$ & $(p\to q)\wedge(p\to r)$ & $q\wedge r$ & $p\to(q\wedge r)$ \\
\midrule
V & V & V & V & V & V & V & V \\
V & V & F & V & F & F & F & F \\
V & F & V & F & V & F & F & F \\
V & F & F & F & F & F & F & F \\
F & V & V & V & V & V & V & V \\
F & V & F & V & V & V & F & V \\
F & F & V & V & V & V & F & V \\
F & F & F & V & V & V & F & V \\
\bottomrule
\end{tabular}
\end{center}
Las columnas sexta y octava coinciden fila a fila, luego las fórmulas son equivalentes.

\noindent\textbf{(b)}
\begin{enumerate}[label=\arabic*.]
  \item Se mueve la negación a través de los cuantificadores: $\neg\forall\varepsilon\,\exists\delta\,\forall x\,(\dots)\equiv\exists\varepsilon\,\forall\delta\,\exists x\,\neg(\dots)$.
  \item Se niega el condicional: $\neg(P\to Q)\equiv P\wedge\neg Q$.
  \item Resultado: $\exists\varepsilon>0\ \forall\delta>0\ \exists x\in\R:\ \abs{x-a}<\delta\ \wedge\ \abs{f(x)-f(a)}\ge\varepsilon$. Las restricciones «$\varepsilon>0$» y «$\delta>0$» no se niegan, porque forman parte del dominio del cuantificador.
\end{enumerate}

\noindent\textbf{(c)}
\begin{enumerate}[label=\arabic*.]
  \item El contrarrecíproco es: si $n$ es par, entonces $n^{2}-6n+5$ es impar.
  \item Sea $n=2m$ con $m\in\Z$. Entonces $n^{2}-6n+5=4m^{2}-12m+5=2(2m^{2}-6m+2)+1$.
  \item Como $2m^{2}-6m+2\in\Z$, el número es impar. Al ser equivalente al contrarrecíproco, el enunciado original queda demostrado.
\end{enumerate}

\noindent\textbf{(d)} Sea $P(n)$ la conjunción de $A(n)\colon F_{2n}=F_n(F_{n-1}+F_{n+1})$ y $B(n)\colon F_{2n+1}=F_n^{2}+F_{n+1}^{2}$.
\begin{enumerate}[label=\arabic*.]
  \item \emph{Caso base} ($n=1$): $F_0=0$, $F_1=1$, $F_2=1$, $F_3=2$. $A(1)$: $F_2=1=F_1(F_0+F_2)=1\cdot(0+1)$. $B(1)$: $F_3=2=F_1^{2}+F_2^{2}=1+1$.
  \item \emph{Paso inductivo}: se supone $P(n)$ para un $n\ge1$. Por la recurrencia y la hipótesis,
  \[
  F_{2n+2}=F_{2n+1}+F_{2n}=F_n^{2}+F_{n+1}^{2}+F_nF_{n-1}+F_nF_{n+1}.
  \]
  Como $F_n^{2}+F_nF_{n-1}=F_n(F_n+F_{n-1})=F_nF_{n+1}$, queda $F_{2n+2}=2F_nF_{n+1}+F_{n+1}^{2}=F_{n+1}(2F_n+F_{n+1})=F_{n+1}(F_n+F_{n+2})$, que es $A(n+1)$.
  \item Usando $A(n+1)$ recién probado y $B(n)$:
  \[
  \begin{aligned}
  F_{2n+3}&=F_{2n+2}+F_{2n+1}=2F_nF_{n+1}+F_{n+1}^{2}+F_n^{2}+F_{n+1}^{2}\\
  &=F_{n+1}^{2}+(F_n+F_{n+1})^{2}=F_{n+1}^{2}+F_{n+2}^{2},
  \end{aligned}
  \]
  que es $B(n+1)$. Por el principio de inducción, $P(n)$ se cumple para todo $n\ge1$.
\end{enumerate}
\resultado (a) equivalentes; (b) $\exists\varepsilon>0\ \forall\delta>0\ \exists x:\ \abs{x-a}<\delta\wedge\abs{f(x)-f(a)}\ge\varepsilon$; (c) y (d) demostrados.
\comprobacion para $n=5$: $F_{10}=55=F_5(F_4+F_6)=5\cdot(3+8)$ y $F_{11}=89=F_5^{2}+F_6^{2}=25+64$. La comprobación numérica hasta $n=24$ también es correcta.

\noindent[Criterios: (a) tabla completa 1~p. (b) cuantificadores 0{,}5~p y condicional 0{,}5~p. (c) enunciar el contrarrecíproco 0{,}5~p, cálculo 1~p, conclusión 0{,}5~p. (d) caso base 0{,}5~p, $A(n+1)$ 1{,}5~p, $B(n+1)$ 1~p; probar una sola identidad sin la otra, como máximo 1~p.]

\noindent Repasar: Módulo 0.1, §1.4, §2.4, §3.3, §5.1 y §5.4.
\end{solucion}

% ------------------------------------------------------------- B.2
\begin{solucion}{eje:funciones}
\datos $f(x)=(2x+3)/(x-1)$ de $\R\setminus\{1\}$ en $\R\setminus\{2\}$; $f\colon A\to B$, $g\colon B\to C$; $f\colon X\to Y$ arbitraria.
\planteamiento se aplican las definiciones de inyectividad y sobreyectividad (Módulo 0.2, §3.1), de imagen y preimagen (Módulo 0.2, §2.3) y los resultados sobre composición (Módulo 0.2, §4.2). En~(c) se demuestra el ciclo (i)$\Rightarrow$(ii)$\Rightarrow$(iii)$\Rightarrow$(i).
\pasos

\noindent\textbf{(a)}
\begin{enumerate}[label=\arabic*.]
  \item \emph{Bien definida}: si $f(x)=2$, entonces $2x+3=2x-2$, es decir, $3=-2$, absurdo. Luego $f(x)\neq2$.
  \item \emph{Inyectiva}: si $f(a)=f(b)$, entonces $(2a+3)(b-1)=(2b+3)(a-1)$, es decir, $2ab-2a+3b-3=2ab-2b+3a-3$, de donde $5b=5a$ y $a=b$.
  \item \emph{Sobreyectiva}: dado $y\neq2$, se busca $x\neq1$ con $f(x)=y$: $y(x-1)=2x+3\iff x(y-2)=y+3\iff x=\frac{y+3}{y-2}$. Este $x$ es distinto de~$1$, pues $\frac{y+3}{y-2}=1$ daría $3=-2$.
  \item \emph{Inversa}: $f^{-1}\colon\R\setminus\{2\}\to\R\setminus\{1\}$, $f^{-1}(y)=\dfrac{y+3}{y-2}$.
\end{enumerate}

\noindent\textbf{(b)}
\begin{enumerate}[label=\arabic*.]
  \item Sea $c\in C$. Como $g\circ f$ es sobreyectiva, existe $a\in A$ con $g(f(a))=c$. Entonces $b=f(a)\in B$ cumple $g(b)=c$, luego $g$ es sobreyectiva.
  \item \emph{Contraejemplo}: $A=\{1\}$, $B=\{1,2\}$, $C=\{1\}$, $f(1)=1$ y $g(1)=g(2)=1$. Entonces $g\circ f\colon\{1\}\to\{1\}$ es sobreyectiva, pero $f$ no lo es, porque $2\notin f(A)$.
\end{enumerate}

\noindent\textbf{(c)} Para todo $A$ se cumple $A\subseteq f^{-1}(f(A))$ y para todos $A,B$ se cumple $f(A\cap B)\subseteq f(A)\cap f(B)$ (Módulo 0.2, §2.3); solo hay que probar los contenidos contrarios.
\begin{enumerate}[label=\arabic*.]
  \item (i)$\Rightarrow$(ii). Si $x\in f^{-1}(f(A))$, existe $a\in A$ con $f(x)=f(a)$; por la inyectividad, $x=a\in A$.
  \item (ii)$\Rightarrow$(iii). Sea $y\in f(A)\cap f(B)$: $y=f(a)=f(b)$ con $a\in A$, $b\in B$. Aplicando (ii) al conjunto $\{a\}$, $b\in f^{-1}(f(\{a\}))=\{a\}$, luego $b=a\in A\cap B$ e $y\in f(A\cap B)$.
  \item (iii)$\Rightarrow$(i). Por reducción al absurdo: si $f(x)=f(x')$ con $x\neq x'$, se toma $A=\{x\}$, $B=\{x'\}$. Entonces $f(A\cap B)=f(\emptyset)=\emptyset$, mientras que $f(A)\cap f(B)=\{f(x)\}\neq\emptyset$, en contra de~(iii).
\end{enumerate}
\resultado (a) $f$ es biyectiva y $f^{-1}(y)=(y+3)/(y-2)$; (b) $g$ es sobreyectiva y $f$ puede no serlo; (c) las tres afirmaciones son equivalentes.
\comprobacion $f(0)=-3$ y $f^{-1}(-3)=0/(-5)=0$; $f(7)=17/6$ y $f^{-1}(17/6)=(35/6)/(5/6)=7$. En~(c), con $f(x)=x^{2}$ en $\R$ (no inyectiva), $A=\{1\}$ da $f^{-1}(f(A))=\{-1,1\}\neq A$, como predice la equivalencia.

\noindent[Criterios: (a) bien definida 0{,}5~p, inyectiva 0{,}5~p, sobreyectiva e inversa 1~p. (b) demostración 1{,}5~p, contraejemplo 0{,}5~p. (c) 1~p por implicación; se descuenta 0{,}5~p si no se mencionan los contenidos que se cumplen siempre.]

\noindent Repasar: Módulo 0.2, §2.3, §3.1, §4.2 y §4.4.
\end{solucion}

% ------------------------------------------------------------- B.3
\begin{solucion}{eje:sumas}
\datos sumandos $1/(k(k+1))$ y $k^{2}$; $x>0$; $V=50\,000$; $\ln5=1{,}6094$, $\ln10=2{,}3026$, $\ln2=0{,}6931$.
\planteamiento (a) descomposición en fracciones simples y suma telescópica (Módulo 0.3, §4.3); (b) inducción (Módulo 0.3, §4.1); (c) un cuadrado es no negativo (Módulo 0.3, §2.4); (d) número de cifras binarias y unidades de información (Módulo 0.3, §5.5 y §6.2).
\pasos

\noindent\textbf{(a)}
\begin{enumerate}[label=\arabic*.]
  \item $\dfrac{1}{k(k+1)}=\dfrac1k-\dfrac{1}{k+1}$, porque $\dfrac{(k+1)-k}{k(k+1)}=\dfrac{1}{k(k+1)}$.
  \item $\displaystyle\sum_{k=1}^{n}\Bigl(\frac1k-\frac1{k+1}\Bigr)=1-\frac{1}{n+1}=\frac{n}{n+1}$, porque los términos intermedios se cancelan.
  \item Como $k^{2}+k=k(k+1)$, las sumas parciales de la serie son $n/(n+1)\to1$. La serie converge y su suma, $1$, es racional.
\end{enumerate}

\noindent\textbf{(b)}
\begin{enumerate}[label=\arabic*.]
  \item \emph{Caso base}: $n=1$ da $1=\frac{1\cdot2\cdot3}{6}$.
  \item \emph{Paso}: si la fórmula vale para $n$, entonces
  \[
  \sum_{k=1}^{n+1}k^{2}=\frac{n(n+1)(2n+1)}{6}+(n+1)^{2}=\frac{(n+1)\bigl(2n^{2}+7n+6\bigr)}{6}=\frac{(n+1)(n+2)(2n+3)}{6},
  \]
  que es la fórmula para $n+1$, pues $2(n+1)+1=2n+3$.
  \item $\displaystyle\sum_{k=11}^{20}k^{2}=\sum_{k=1}^{20}k^{2}-\sum_{k=1}^{10}k^{2}=\frac{20\cdot21\cdot41}{6}-\frac{10\cdot11\cdot21}{6}=2870-385=2485$.
\end{enumerate}

\noindent\textbf{(c)} Para $x>0$, $x+\frac1x-2=\frac{x^{2}-2x+1}{x}=\frac{(x-1)^{2}}{x}\ge0$, y vale $0$ si y solo si $x=1$. Es la desigualdad entre las medias aritmética y geométrica de $x$ y $1/x$, cuyo producto es~$1$.

\noindent\textbf{(d)}
\begin{enumerate}[label=\arabic*.]
  \item Con $b$ bits se numeran $2^{b}$ valores; hace falta $2^{b}\ge50\,000$. Como $2^{15}=32\,768<50\,000\le65\,536=2^{16}$, basta con $b=16$ bits.
  \item $\ln50\,000=\ln5+4\ln10=1{,}6094+9{,}2104=10{,}82$ nats.
  \item $\log_2 50\,000=\ln50\,000/\ln2=10{,}82/0{,}6931=15{,}61$ bits.
\end{enumerate}
\resultado (a) $n/(n+1)$; la serie suma $1$, racional. (b) $2485$. (c) Demostrado. (d) 16 bits; $15{,}61$ bits $=10{,}82$ nats.
\comprobacion con $n=3$, $\frac12+\frac16+\frac1{12}=\frac34$. La suma directa de $121+144+\dots+400$ da $2485$. Además, $15{,}61\le16$: el código de longitud fija nunca usa menos bits que la información de un símbolo uniforme.

\noindent[Criterios: (a) descomposición 0{,}5~p, telescópica 1~p, serie 0{,}5~p. (b) base 0{,}5~p, paso 1~p, cálculo 0{,}5~p. (c) 1~p, de los que 0{,}5~p corresponden al caso de igualdad. (d) bits del identificador 0{,}5~p e información 0{,}5~p.]

\noindent Repasar: Módulo 0.3, §2.4, §4.1, §4.3, §5.5 y §6.2.
\end{solucion}

% ------------------------------------------------------------- B.4
\begin{solucion}{eje:complejos}
\datos $\theta\notin\frac{\pi}{2}\Z$ en (a); $w=8i$; el cociente $(1+i)^{2025}/(1-i)^{2023}$.
\planteamiento (a) fórmula del seno de una diferencia y del ángulo doble (Módulo 0.4, §2.2 y §2.3); (b) y (c) teorema de De Moivre y raíces $n$-ésimas (Módulo 0.4, §6.4 y §6.5); (d) cociente en forma binómica o polar (Módulo 0.4, §6.2).
\pasos

\noindent\textbf{(a)}
\begin{enumerate}[label=\arabic*.]
  \item Si $\theta$ no es múltiplo de $\pi/2$, $\sin\theta\neq0$ y $\cos\theta\neq0$, de modo que la expresión está definida.
  \item Con denominador común: $\dfrac{\sin3\theta\cos\theta-\cos3\theta\sin\theta}{\sin\theta\cos\theta}=\dfrac{\sin(3\theta-\theta)}{\sin\theta\cos\theta}$.
  \item $\sin2\theta=2\sin\theta\cos\theta$, luego el cociente vale $2$.
\end{enumerate}

\noindent\textbf{(b)} Sean $c=\cos\theta$ y $s=\sin\theta$.
\begin{enumerate}[label=\arabic*.]
  \item Por De Moivre, $\cos3\theta+i\sin3\theta=(c+is)^{3}$.
  \item Por el binomio, $(c+is)^{3}=c^{3}+3c^{2}(is)+3c(is)^{2}+(is)^{3}=(c^{3}-3cs^{2})+i(3c^{2}s-s^{3})$.
  \item Igualando partes reales y usando $s^{2}=1-c^{2}$: $\cos3\theta=c^{3}-3c(1-c^{2})=4c^{3}-3c$.
  \item Igualando partes imaginarias y usando $c^{2}=1-s^{2}$: $\sin3\theta=3(1-s^{2})s-s^{3}=3s-4s^{3}$.
\end{enumerate}

\noindent\textbf{(c)}
\begin{enumerate}[label=\arabic*.]
  \item \emph{De Moivre}: para todo $\theta\in\R$ y todo $n\in\Z$, $(\cos\theta+i\sin\theta)^{n}=\cos n\theta+i\sin n\theta$, es decir, $(e^{i\theta})^{n}=e^{in\theta}$.
  \item Se busca $z=re^{i\theta}$ con $z^{n}=w=\rho e^{i\varphi}$. Por De Moivre, $r^{n}e^{in\theta}=\rho e^{i\varphi}$, lo que equivale a $r^{n}=\rho$ y $n\theta=\varphi+2k\pi$ con $k\in\Z$.
  \item Luego $z_k=\rho^{1/n}e^{i(\varphi+2k\pi)/n}$; los valores $k=0,\dots,n-1$ dan $n$ raíces distintas, y los demás $k$ las repiten.
  \item Para $w=8i=8e^{i\pi/2}$ y $n=3$: $z_k=2e^{i(\pi/6+2k\pi/3)}$. Así, $z_0=2\bigl(\frac{\sqrt3}{2}+\frac i2\bigr)=\sqrt3+i$, $z_1=2e^{i5\pi/6}=-\sqrt3+i$ y $z_2=2e^{i3\pi/2}=-2i$.
  \item Suma: $(\sqrt3+i)+(-\sqrt3+i)+(-2i)=0$.
\end{enumerate}

\noindent\textbf{(d)}
\begin{enumerate}[label=\arabic*.]
  \item $\dfrac{1+i}{1-i}=\dfrac{(1+i)^{2}}{(1-i)(1+i)}=\dfrac{2i}{2}=i$.
  \item $\dfrac{(1+i)^{2025}}{(1-i)^{2023}}=\Bigl(\dfrac{1+i}{1-i}\Bigr)^{2023}(1+i)^{2}=i^{2023}\cdot2i=2\,i^{2024}$.
  \item Como $2024=4\cdot506$, $i^{2024}=1$ y el cociente vale~$2$.
\end{enumerate}
\resultado (a) y (b) demostrados; (c) $\sqrt3+i$, $-\sqrt3+i$, $-2i$; (d) $2$.
\comprobacion $(\sqrt3+i)^{2}=2+2\sqrt3\,i$ y $(2+2\sqrt3\,i)(\sqrt3+i)=2\sqrt3-2\sqrt3+(2+6)i=8i$. En forma polar, en~(d): $1\pm i=\sqrt2e^{\pm i\pi/4}$, luego el cociente es $(\sqrt2)^{2}e^{i(2025+2023)\pi/4}=2e^{i\,1012\pi}=2$. Con $\theta=0{,}37$, la expresión de~(a) vale $2{,}0000$.

\noindent[Criterios: (a) dominio 0{,}5~p, diferencia 1~p, ángulo doble 0{,}5~p. (b) desarrollo 1~p, cada identidad 0{,}5~p. (c) enunciado 0{,}5~p, deducción 0{,}5~p, raíces 1~p. (d) 1~p.]

\noindent Repasar: Módulo 0.4, §2.2, §2.3, §6.2, §6.4 y §6.5.
\end{solucion}

% ------------------------------------------------------------- B.5
\begin{solucion}{eje:tvm}
\datos $\lim_{x\to2}(x^{2}+x)$; $f(x)=x/(x+1)$; $\arctan$ en $[0,x]$, con $\arctan'(t)=1/(1+t^{2})$.
\planteamiento (a) se acota $\abs{x+3}$ en un entorno de~$2$ (Módulo 0.5, §2.3); (b) cociente incremental (Módulo 0.5, §4.1); (c) teorema del valor medio de Lagrange en $[0,x]$ (Módulo 0.5, §5.2); (d) regla de L'Hôpital o polinomio de Taylor (Módulo 0.5, §5.3 y §6.4).
\pasos

\noindent\textbf{(a)}
\begin{enumerate}[label=\arabic*.]
  \item $\abs{x^{2}+x-6}=\abs{x-2}\,\abs{x+3}$.
  \item Si $\abs{x-2}<1$, entonces $1<x<3$ y $4<x+3<6$, luego $\abs{x+3}<6$.
  \item Dado $\varepsilon>0$, se toma $\delta=\min\{1,\varepsilon/6\}$. Si $0<\abs{x-2}<\delta$, entonces $\abs{x^{2}+x-6}<6\delta\le\varepsilon$.
\end{enumerate}

\noindent\textbf{(b)} Para $x\neq-1$ y $h\neq0$ con $x+h\neq-1$:
\[
\begin{aligned}
\frac{f(x+h)-f(x)}{h}&=\frac{(x+h)(x+1)-x(x+h+1)}{h\,(x+h+1)(x+1)}=\frac{h}{h\,(x+h+1)(x+1)}\\
&=\frac{1}{(x+h+1)(x+1)}\xrightarrow[h\to0]{}\frac{1}{(x+1)^{2}}.
\end{aligned}
\]

\noindent\textbf{(c)}
\begin{enumerate}[label=\arabic*.]
  \item Sea $x>0$. La función $\arctan$ es continua en $[0,x]$ y derivable en $(0,x)$.
  \item Por el teorema del valor medio, existe $c\in(0,x)$ tal que $\arctan x-\arctan0=\arctan'(c)\,(x-0)$, es decir, $\arctan x=\dfrac{x}{1+c^{2}}$.
  \item De $0<c<x$ se sigue $1<1+c^{2}<1+x^{2}$, y por tanto $\dfrac{1}{1+x^{2}}<\dfrac{1}{1+c^{2}}<1$.
  \item Al multiplicar por $x>0$, las desigualdades se conservan: $\dfrac{x}{1+x^{2}}<\arctan x<x$.
\end{enumerate}

\noindent\textbf{(d)} Es una indeterminación $0/0$. Por L'Hôpital,
\[
\lim_{x\to0}\frac{1-\frac{1}{1+x^{2}}}{3x^{2}}=\lim_{x\to0}\frac{x^{2}}{3x^{2}(1+x^{2})}=\lim_{x\to0}\frac{1}{3(1+x^{2})}=\frac13 .
\]
También puede usarse $\arctan x=x-\frac{x^{3}}{3}+o(x^{3})$.
\resultado (a) demostrado con $\delta=\min\{1,\varepsilon/6\}$; (b) $f'(x)=1/(x+1)^{2}$; (c) demostrado; (d) $1/3$.
\comprobacion con $x=1$: $0{,}5<\arctan1=0{,}7854<1$. Con $x=0{,}01$, el cociente de~(d) vale $0{,}33331$. En~(a), para $\varepsilon=0{,}6$ se obtiene $\delta=0{,}1$, y con $x=2{,}09$ resulta $\abs{x^{2}+x-6}=0{,}4581<0{,}6$.

\noindent[Criterios: (a) factorización 0{,}5~p, acotación 0{,}5~p, elección de $\delta$ y cierre 1~p. (b) 1~p. (c) hipótesis del teorema 0{,}5~p, aplicación 1~p, desigualdades 1{,}5~p; sin justificar el signo de~$x$ al multiplicar, $-0{,}5$~p. (d) 1~p.]

\noindent Repasar: Módulo 0.5, §2.3, §4.1, §5.2, §5.3 y §6.4.
\end{solucion}

% ------------------------------------------------------------- B.6
\begin{solucion}{eje:integrales}
\datos $\int_1^{e}x\ln x\,\dd x$; $\int_1^{\infty}\ln x/x^{2}\,\dd x$; $\ln(1+x)$ en torno a~$0$; $x=0{,}2$; $e^{2}=7{,}3891$.
\planteamiento integración por partes (Módulo 0.6, §3.2), integral impropia de primera especie (Módulo 0.6, §4.1), integración término a término de la serie geométrica (Módulo 0.6, §6.2 y §6.4), radio de convergencia (Módulo 0.6, §6.1) y resto de Lagrange (Módulo 0.6, §6.3; Módulo 0.5, §6.2).
\pasos

\noindent\textbf{(a)} Con $u=\ln x$ y $\dd v=x\,\dd x$ ($\dd u=\dd x/x$, $v=x^{2}/2$):
\[
\int_1^{e}x\ln x\,\dd x=\Bigl[\frac{x^{2}}{2}\ln x\Bigr]_1^{e}-\int_1^{e}\frac{x}{2}\,\dd x=\frac{e^{2}}{2}-\Bigl[\frac{x^{2}}{4}\Bigr]_1^{e}=\frac{e^{2}}{2}-\frac{e^{2}}{4}+\frac14=\frac{e^{2}+1}{4}=2{,}097.
\]

\noindent\textbf{(b)}
\begin{enumerate}[label=\arabic*.]
  \item Para $b>1$, por partes con $u=\ln x$ y $\dd v=x^{-2}\dd x$ ($v=-1/x$): $\displaystyle\int_1^{b}\frac{\ln x}{x^{2}}\,\dd x=\Bigl[-\frac{\ln x}{x}\Bigr]_1^{b}+\int_1^{b}\frac{\dd x}{x^{2}}=-\frac{\ln b}{b}+1-\frac1b$.
  \item Como $\ln b/b\to0$ cuando $b\to\infty$ (el logaritmo crece más despacio que cualquier potencia), el límite existe: la integral converge y vale~$1$.
\end{enumerate}

\noindent\textbf{(c)}
\begin{enumerate}[label=\arabic*.]
  \item Para $\abs{t}<1$, $\dfrac{1}{1+t}=\sum_{n\ge0}(-1)^{n}t^{n}$ (serie geométrica de razón $-t$).
  \item Integrando término a término entre $0$ y $x$, con $\abs x<1$: $\ln(1+x)=\displaystyle\sum_{n\ge1}\frac{(-1)^{n+1}}{n}x^{n}=x-\frac{x^{2}}{2}+\frac{x^{3}}{3}-\cdots$.
  \item Radio: $\abs{a_{n+1}/a_n}=n/(n+1)\to1$, luego $R=1$.
  \item $\ln1{,}2\approx0{,}2-\frac{0{,}2^{2}}{2}=0{,}2-0{,}02=0{,}18$.
\end{enumerate}

\noindent\textbf{(d)}
\begin{enumerate}[label=\arabic*.]
  \item La aproximación es $P_2(x)=x-x^{2}/2$. El resto es $R_2(x)=\dfrac{f'''(c)}{3!}x^{3}$ con $c\in(0;\,0{,}2)$ y $f'''(t)=\dfrac{2}{(1+t)^{3}}$.
  \item Para $0<c<0{,}2$, $0<f'''(c)<2$; por tanto, $0<R_2(0{,}2)<\dfrac{2}{6}\cdot0{,}008=0{,}002667$.
  \item Error real: $0{,}18232-0{,}18=0{,}00232$, positivo y menor que la cota.
\end{enumerate}
\resultado (a) $(e^{2}+1)/4=2{,}097$; (b) converge a~$1$; (c) $\ln(1+x)=\sum_{n\ge1}(-1)^{n+1}x^{n}/n$, $R=1$, $\ln1{,}2\approx0{,}18$; (d) error menor que $0{,}002667$.
\comprobacion la integración numérica da $2{,}0973$ y $1{,}0000$. En~(d), la cota coincide con el tercer término de la serie alternada, $0{,}2^{3}/3=0{,}002667$, que también acota el error por el criterio de Leibniz.

\noindent[Criterios: (a) elección de partes 0{,}5~p, cálculo 1~p. (b) primitiva 0{,}5~p, límite justificado 1~p. (c) serie 1~p, radio 0{,}5~p, aproximación 0{,}5~p. (d) 1~p.]

\noindent Repasar: Módulo 0.6, §3.2, §4.1, §6.1, §6.2, §6.3 y §6.4; Módulo 0.5, §5.3 (jerarquía de crecimientos) y §6.2.
\end{solucion}

% ------------------------------------------------------------- B.7
\begin{solucion}{eje:varias}
\datos $f(x,y)=x^{2}y^{3}-2x$ en $(2,1)$; $g(x,y)=x^{3}-3x+y^{2}-4y$; $h(x,y)=xy$ con $x^{2}+4y^{2}=8$; $D=\{x^{2}+y^{2}\le4\}$.
\planteamiento derivadas parciales y plano tangente (Módulo 0.7, §4.1 y §4.2), derivada direccional y gradiente (Módulo 0.7, §4.4), criterio del hessiano (Módulo 0.7, §5.2), Weierstrass y multiplicadores de Lagrange (Módulo 0.7, §5.3 y §5.4) y coordenadas polares (Módulo 0.7, §6.3).
\pasos

\noindent\textbf{(a)}
\begin{enumerate}[label=\arabic*.]
  \item $f(2,1)=4-4=0$; $f_x=2xy^{3}-2$ y $f_y=3x^{2}y^{2}$, luego $\nabla f(2,1)=(2,12)$.
  \item Plano tangente: $z=0+2(x-2)+12(y-1)$, es decir, $z=2x+12y-16$.
\end{enumerate}

\noindent\textbf{(b)}
\begin{enumerate}[label=\arabic*.]
  \item $f(2{,}1;0{,}9)\approx2\cdot0{,}1+12\cdot(-0{,}1)=-1{,}0$. Valor exacto: $4{,}41\cdot0{,}729-4{,}2=-0{,}9851$; el error es $0{,}0149$.
  \item Con $u=(3,4)/5$: $D_uf(2,1)=\nabla f\cdot u=(6+48)/5=10{,}8$.
  \item El máximo de la derivada direccional es $\norm{\nabla f(2,1)}=\sqrt{4+144}=\sqrt{148}=12{,}17$, en la dirección de $(2,12)$.
\end{enumerate}

\noindent\textbf{(c)}
\begin{enumerate}[label=\arabic*.]
  \item $g_x=3x^{2}-3=0\iff x=\pm1$; $g_y=2y-4=0\iff y=2$. Puntos críticos: $(1,2)$ y $(-1,2)$.
  \item Hessiana: $H=\begin{pmatrix}6x&0\\0&2\end{pmatrix}$.
  \item En $(1,2)$: $\det H=12>0$ y $g_{xx}=6>0$: mínimo local, con $g(1,2)=1-3+4-8=-6$.
  \item En $(-1,2)$: $\det H=-12<0$: punto de silla, con $g(-1,2)=-2$.
\end{enumerate}

\noindent\textbf{(d)}
\begin{enumerate}[label=\arabic*.]
  \item La elipse es cerrada y acotada y $h$ es continua: por el teorema de Weierstrass, $h$ alcanza máximo y mínimo en ella. Además, $\nabla(x^{2}+4y^{2})=(2x,8y)\neq(0,0)$ en la elipse, así que el método de Lagrange es aplicable.
  \item Ecuaciones: $y=2\lambda x$, $x=8\lambda y$, $x^{2}+4y^{2}=8$. Sustituyendo, $x=16\lambda^{2}x$, luego $x=0$ o $\lambda=\pm\frac14$. Si $x=0$, entonces $y=0$, que no está en la elipse.
  \item $\lambda=\frac14$: $y=x/2$, luego $2x^{2}=8$ y se obtienen $(2,1)$ y $(-2,-1)$, con $h=2$.
  \item $\lambda=-\frac14$: $y=-x/2$, luego se obtienen $(2,-1)$ y $(-2,1)$, con $h=-2$.
\end{enumerate}

\noindent\textbf{(e)} En polares, $\dd A=r\,\dd r\,\dd\theta$:
\[
\int_0^{2\pi}\!\!\int_0^{2}e^{-r^{2}}r\,\dd r\,\dd\theta=2\pi\Bigl[-\tfrac12e^{-r^{2}}\Bigr]_0^{2}=\pi\bigl(1-e^{-4}\bigr)=3{,}1416\cdot0{,}9817=3{,}084.
\]
\resultado (a) $\nabla f(2,1)=(2,12)$, $z=2x+12y-16$; (b) $-1{,}0$ frente a $-0{,}9851$, $D_uf=10{,}8$, máximo $12{,}17$; (c) mínimo local en $(1,2)$ y silla en $(-1,2)$; (d) máximo $2$ en $(2,1)$ y $(-2,-1)$, mínimo $-2$ en $(2,-1)$ y $(-2,1)$; (e) $\pi(1-e^{-4})=3{,}084$.
\comprobacion en~(d), $2^{2}+4\cdot1^{2}=8$; además, por la desigualdad de las medias, $\abs{xy}=\frac12\abs{x}\,\abs{2y}\le\frac14(x^{2}+4y^{2})=2$, que confirma los valores extremos. En~(e), la integración numérica da $3{,}0841$.

\noindent[Criterios: (a) gradiente 0{,}5~p, plano 0{,}5~p. (b) aproximación 0{,}5~p, derivadas direccionales 0{,}5~p. (c) puntos 1~p, clasificación 1~p. (d) existencia 0{,}5~p, sistema y casos 1~p, valores 0{,}5~p. (e) 1~p.]

\noindent Repasar: Módulo 0.7, §4.1, §4.2, §4.4, §5.2, §5.3, §5.4 y §6.3.
\end{solucion}

% ------------------------------------------------------------- B.8
\begin{solucion}{eje:probabilidad}
\datos vocabulario de 12 tokens; monedas con $p_1=0{,}5$, $p_2=0{,}7$, $p_3=0{,}2$, elegidas con probabilidad $1/3$ cada una; tres lanzamientos independientes dada la moneda.
\planteamiento (a) variaciones y combinaciones (Módulo 0.8, §1.4 y §2.1); (b) doble conteo (Módulo 0.8, §2.4); (c) probabilidad total y Bayes (Módulo 0.8, §4.3); (d) esperanza condicionada a cada moneda, linealidad y binomial (Módulo 0.8, §5.2 y §5.4).
\pasos

\noindent\textbf{(a)} Con orden: $V_{12,4}=12\cdot11\cdot10\cdot9=11\,880$. Sin orden: $\binom{12}{4}=11\,880/4!=495$.

\noindent\textbf{(b)}
\begin{enumerate}[label=\arabic*.]
  \item Se cuentan los subconjuntos de $k$ elementos de $\{1,\dots,n\}$: hay $\binom nk$.
  \item Se clasifican según contengan o no el elemento~$n$. Los que lo contienen quedan determinados por los otros $k-1$ elementos, elegidos entre $n-1$: hay $\binom{n-1}{k-1}$. Los que no lo contienen son subconjuntos de $k$ elementos de $\{1,\dots,n-1\}$: hay $\binom{n-1}{k}$.
  \item Las dos clases son disjuntas y cubren todos los casos; por la regla de la suma, $\binom nk=\binom{n-1}{k-1}+\binom{n-1}{k}$.
\end{enumerate}

\noindent\textbf{(c)} Sea $S$ la secuencia cara, cruz, cruz y $M_j$ el suceso «se elige la moneda~$j$».
\begin{enumerate}[label=\arabic*.]
  \item $\PP(S\mid M_j)=p_j(1-p_j)^{2}$: $0{,}5\cdot0{,}25=0{,}125$; $0{,}7\cdot0{,}09=0{,}063$; $0{,}2\cdot0{,}64=0{,}128$.
  \item Probabilidad total: $\PP(S)=\frac13(0{,}125+0{,}063+0{,}128)=0{,}316/3=0{,}1053$.
  \item Bayes: $\PP(M_1\mid S)=\dfrac{\frac13\cdot0{,}125}{\frac13\cdot0{,}316}=0{,}3956$.
\end{enumerate}

\noindent\textbf{(d)}
\begin{enumerate}[label=\arabic*.]
  \item Dada $M_j$, $X\sim\Bin(3,p_j)$, con $\mathbb E[X\mid M_j]=3p_j$. Por probabilidad total, $\mathbb E[X]=\frac13(1{,}5+2{,}1+0{,}6)=1{,}4$. Equivalentemente, por linealidad: cada lanzamiento es cara con probabilidad $\frac13(0{,}5+0{,}7+0{,}2)=0{,}4667$, y $3\cdot0{,}4667=1{,}4$.
  \item $\PP(X=3)=\frac13\bigl(0{,}5^{3}+0{,}7^{3}+0{,}2^{3}\bigr)=\frac13(0{,}125+0{,}343+0{,}008)=0{,}476/3=0{,}1587$.
\end{enumerate}
\resultado (a) $11\,880$ y $495$; (b) demostrado; (c) $\PP(S)=0{,}1053$, $\PP(M_1\mid S)=0{,}3956$; (d) $\mathbb E[X]=1{,}4$, $\PP(X=3)=0{,}1587$.
\comprobacion las probabilidades a posteriori suman~1: $(0{,}125+0{,}063+0{,}128)/0{,}316=1$. Si los lanzamientos fueran independientes con probabilidad $0{,}4667$, se obtendría $\PP(X=3)=0{,}4667^{3}=0{,}1016\neq0{,}1587$: dada la mezcla de monedas, los lanzamientos no son independientes, aunque la esperanza sí se calcula por linealidad.

\noindent[Criterios: (a) 0{,}5~p cada valor. (b) 1~p. (c) verosimilitudes 0{,}5~p, total 0{,}5~p, Bayes 1~p. (d) esperanza 1~p, $\PP(X=3)$ 1~p.]

\noindent Repasar: Módulo 0.8, §1.4, §2.1, §2.4, §4.3, §5.2 y §5.4.
\end{solucion}

% ------------------------------------------------------------- B.9
\begin{solucion}{eje:estadistica}
\datos muestra de Pareto $2,4,8$; $n=16$, $\sigma=4$, $\xbar=3{,}9$, $\mu_0=2$, $\alpha=0{,}05$; latencias $12,15,11,14,13,16,12,11$; 20 pruebas al nivel $0{,}05$.
\planteamiento (a) máxima verosimilitud (Módulo 0.9, §4.2); (b) contraste $z$ bilateral y $p$-valor (Módulo 0.9, §5.1 a §5.3); (c) intervalo $t$ (Módulo 0.9, §1.4 y §4.4); (d) tasa de error por familia y Bonferroni (Módulo 0.9, §5.6).
\pasos

\noindent\textbf{(a)}
\begin{enumerate}[label=\arabic*.]
  \item Verosimilitud: $L(\alpha)=\prod_{i=1}^{3}\alpha x_i^{-\alpha-1}=\alpha^{3}\,(x_1x_2x_3)^{-\alpha-1}$.
  \item Log-verosimilitud: $\ell(\alpha)=3\ln\alpha-(\alpha+1)\sum\ln x_i$, con $\sum\ln x_i=\ln2+\ln4+\ln8=6\ln2=4{,}159$.
  \item $\ell'(\alpha)=\dfrac3\alpha-6\ln2=0\iff\widehat\alpha=\dfrac{1}{2\ln2}=0{,}7213$.
  \item $\ell''(\alpha)=-3/\alpha^{2}<0$ para todo $\alpha>0$: $\ell$ es cóncava y el punto crítico es el máximo global.
\end{enumerate}

\noindent\textbf{(b)}
\begin{enumerate}[label=\arabic*.]
  \item Bajo $H_0$, $\Xbar\sim\Nor\bigl(2,(4/\sqrt{16})^{2}\bigr)=\Nor(2,1)$.
  \item Región de rechazo: $\abs{\xbar-2}>1{,}96$, es decir, $\xbar<0{,}04$ o $\xbar>3{,}96$.
  \item $z=(3{,}9-2)/1=1{,}90$; $p=2\bigl(1-\Phi(1{,}90)\bigr)=2\cdot0{,}0287=0{,}0574$.
  \item Como $p=0{,}0574>0{,}05$ (equivalentemente, $3{,}9<3{,}96$), no se rechaza $H_0$ al nivel del 5\,\%.
\end{enumerate}

\noindent\textbf{(c)}
\begin{enumerate}[label=\arabic*.]
  \item $\xbar=104/8=13$.
  \item Desviaciones: $-1,2,-2,1,0,3,-1,-2$; suma de cuadrados $24$; $s^{2}=24/7=3{,}429$ y $s=1{,}852$.
  \item Con $n-1=7$ grados de libertad, $t_{7;0{,}975}=2{,}365$; el margen es $2{,}365\cdot1{,}852/\sqrt8=1{,}548$.
  \item Intervalo: $13\pm1{,}548=[11{,}45;\ 14{,}55]$~ms.
\end{enumerate}

\noindent\textbf{(d)} Por independencia, $\PP(\text{ningún falso positivo})=0{,}95^{20}=0{,}3585$, luego $\PP(\text{al menos uno})=0{,}6415$. Bonferroni: cada prueba al nivel $0{,}05/20=0{,}0025$; por la cota de la unión, la probabilidad de algún falso positivo es como máximo $20\cdot0{,}0025=0{,}05$.
\resultado (a) $\widehat\alpha=1/(2\ln2)=0{,}7213$; (b) $z=1{,}90$, $p=0{,}0574$, no se rechaza $H_0$; (c) $\xbar=13$, $s=1{,}852$, IC $[11{,}45;\ 14{,}55]$; (d) $0{,}6415$ y nivel $0{,}0025$.
\comprobacion en~(a), $\ell(0{,}7)=-8{,}1401$, $\ell(0{,}7213)=-8{,}1388$ y $\ell(0{,}75)=-8{,}1411$: el valor central es el mayor. En~(c), usar $t_{8;0{,}975}=2{,}306$ (grados de libertad incorrectos) daría un margen de $1{,}510$; usar $z_{0{,}975}=1{,}96$, de $1{,}283$.

\noindent[Criterios: (a) log-verosimilitud 1~p, estimador 0{,}5~p, segunda derivada 0{,}5~p. (b) región 0{,}5~p, $z$ 0{,}5~p, $p$-valor 0{,}5~p, decisión 0{,}5~p. (c) media y $s$ 1~p, cuantil con 7 grados de libertad 0{,}5~p, intervalo 0{,}5~p. (d) 0{,}5~p cada parte.]

\noindent Repasar: Módulo 0.9, §1.4, §4.2, §4.4, §5.1, §5.2, §5.3 y §5.6.
\end{solucion}
